\documentclass{article}

\usepackage[margin=0.75in]{geometry}
\usepackage{amsmath}
\usepackage{amsfonts}
\usepackage{multicol}
\usepackage{hyperref}
\usepackage{titling}

\setlength{\parindent}{0pt}
\setlength{\droptitle}{-1cm}

\title{Stochastic Processes Notes}
\author{gordonbchen}
\begin{document}
\date{}
\maketitle


\section{Chapter 1}
\subsection{Markov Property}
Markov property: $\boxed{\mathbb{P}(X_{n+1} | X_n, X_{n-1} \dots X_1, X_0) = \mathbb{P}(X_{n+1} | X_n)}$. \\

$P_{i,j} = \mathbb{P}(X_{n+1}=j | X_n=i)$. Since $\sum_j \mathbb{P}(X_{n+1}=j | X_n=i) = 1$, then \boxed{P \vec{1} = \vec{1}}. \\

$m$-step transition probability: \boxed{\mathbb{P}(X_{n+m}=j | X_n=i) = P^m(i,j)}.


\subsection{Classifications}
\subsubsection{Transient and Recurrent States}
Let $\mathbb{P}_y(T_y < \infty) = \rho_{yy}$ be the probability of returning to $y$, where $T_y = \min\{n \ge 1: X_n = y\}$ is the time of first return to $y$. Then the probability of returning to $y$, $k$ times is $\mathbb{P}_y(T^k_y < \infty) = \rho_{yy}^k$. \\

If $\rho_{yy} < 1$, then $\lim_{k\to\infty} \rho^k_{yy} = 0$ and the state is transient. The probability for visiting the state more and more times goes to 0, and at some point, the state is never visited again. \\

Otherwise, if $\rho_{yy} = 1$, then $\lim_{k\to\infty} \rho^k_{yy} = 1$ and the state is recurrent. The state will be visited infinitely many times. \\

If a set is finite, closed (cannot leave set), and irreducible (all states communicate with each other), then all of its states are recurrent.

\subsubsection{Period}
Period of Markov Chain: greatest common divisor of all $I_x = \{n \ge 1: p^n(x, x) > 0\}$, length of all cycles. \\

Aperiodic Markov Chain: period = 1.


\subsection{Stationary Distribution}
Stationary distribution $\pi$: \boxed{\pi P = \pi}

\subsubsection{Spectral Theory}
$A = \begin{pmatrix}
a & 1-a \\ 
1-b & b
\end{pmatrix}$.
$A$ has eigenvalues $\lambda_1 = 1, \lambda_2 = a+b-1$. $|\lambda_2| < 1$.\\

$\vec{1}$ is a right eigenvalue of $A$, $A \vec{1} = \vec{1}$. $\pi$ is a left eigenvalue of $A$, $\pi A = \pi$. \\

$A = P D P^{-1}$, where $P$ are the right eigenvectors of $A$, $D$ is a diagonal matrix of eigenvalues, and $P^{-1}$ is the matrix of left eigenvectors (rows). To see this last part, note that $A^T = (P^{-1})^T D P^T$ and left eigenvectors are eigenvectors of $A^T$. \\

$A = \begin{pmatrix}\vec{r}_1 & \vec{r}_2\end{pmatrix} \begin{pmatrix}\lambda_1 & 0 \\ 0 & \lambda_2\end{pmatrix} \begin{pmatrix}\vec{l}_1 \\ \vec{l}_2\end{pmatrix}$
$= \lambda_1 \vec{r}_1 \vec{l}_1 + \lambda_2 \vec{r}_2 \vec{l}_2$
$=\begin{pmatrix} \pi_1 & \pi_2 \\ \pi_1 & \pi_2 \end{pmatrix} + \lambda_2 \vec{r_2}\vec{l_2}$ \\

$\lim_{n\to\infty} A^n = P D^n P^{-1} = \begin{pmatrix} \pi_1 & \pi_2 \\ \pi_1 & \pi_2 \end{pmatrix} + \lambda_2^n \vec{r_2}\vec{l_2}$
$= \begin{pmatrix} \pi_1 & \pi_2 \\ \pi_1 & \pi_2 \end{pmatrix}$
since $|\lambda_2| < 1$
\subsection{Limit Distributions}


\subsection{Exit Distributions}
$V_F = \min\{n \ge 0: X_n \in F\}$ is the Hitting Time of set $F$, the first time $X$ takes on a state in $F$. \\

$\mathbb{P}_x(V_A < V_B)$ is the probability that $X$, starting in state $X_0 = x$, takes on a state in $A$ (wins) before taking on a state in $B$ (loses). \\

Let $\mathbb{P}_x(V_A < V_B) = h(x)$. $h(a) = 1$ for all $a \in A$, and $h(b) = 0$ for all $b \in B$. \boxed{h(x) = \sum_y P_{x,y} h(y)}. \\

\subsubsection{Example}\label{subsec:fair_gambler}
A gambler wins \$1 with $p=1/2$ and loses \$1 with $p=1/2$. If he ever reaches \$0, the casino makes him stop. If he ever reaches \$4, he is happy and goes home. Let $X \in \{0, 1, 2, 3, 4\}$ be the amount of money the gambler has. Note that $X = 0$ and $X = 4$ are absorbing states. Let $A = \{4\}$ and $B = \{0\}$, so $h(x) = \mathbb{P}_x(V_A < V_B)$ is the probability that the gambler ends up going home with \$4 given that he currently has $\$x$.


\noindent
\begin{minipage}[t]{0.2\textwidth}
\begin{align*}
h(0) &= 0 \\
h(1) &= \tfrac{1}{2} h(2) \\
h(2) &= \tfrac{1}{2} h(1) + \tfrac{1}{2} h(3) \\
h(3) &= \tfrac{1}{2} h(2) + \tfrac{1}{2} \\
h(4) &= 1
\end{align*}
\end{minipage}
\hspace{1.5em}
\begin{minipage}[t]{0.40\textwidth}
$$\left(
\begin{array}{c|ccc|c}
1   & 0   & 0   & 0   & 0 \\
\hline
0  & 1  & -.5  & 0   & 0 \\
0   & -.5  & 1  & -.5  & 0 \\
0   & 0   & -.5  & 1  & 0 \\
\hline
0   & 0   & 0   & 0   & 1
\end{array}
\right)
\begin{pmatrix}
h(0) \\
h(1) \\
h(2) \\
h(3) \\
h(4)
\end{pmatrix}
=
\begin{pmatrix}
0 \\ 0 \\ 0 \\ .5 \\ 1
\end{pmatrix}$$
\end{minipage}
\hspace{2em}
\begin{minipage}[t]{0.24\textwidth}
\vspace*{1\baselineskip}
$$\begin{pmatrix}
1  & -.5  & 0  \\
-.5  & 1  & -.5 \\
0   & -.5  & 1 \\
\end{pmatrix}
h
=
\begin{pmatrix}
0 \\ 0 \\ .5
\end{pmatrix}$$
\end{minipage} \\ \\

Let $C = S-A-B$, $R = P_{x \in C, y \in C}$, and $v_i = P_{i,a}$. \\

\boxed{h = R h + v}. then $(I-R) h = v$ and $h = (I-R)^{-1} v$.


\subsubsection{Limit distributions with transient states}
\noindent
\begin{minipage}[t]{0.5\textwidth}
$$P = \begin{array}{c|cc|cc|ccc}
& 1 & 2 & 3 & 4 & 5 & 6 & 7 \\
\hline
1 & .2 & .3 & .1 & 0 & .4 & 0 & 0 \\
2 & 0 & .5 & 0 & .2 & .3 & 0 & 0 \\
\hline
3 & 0 & 0 & .7 & .3 & 0 & 0 & 0 \\
4 & 0 & 0 & .6 & .4 & 0 & 0 & 0 \\
\hline
5 & 0 & 0 & 0 & 0 & .5 & .5 & 0 \\
6 & 0 & 0 & 0 & 0 & 0 & .2 & .8 \\
6 & 0 & 0 & 0 & 0 & 1 & 0 & 0 \\
\end{array}$$

$\{1, 2\}$ are transient states. $A = \{3, 4\}$ and $B = \{5, 6, 7\}$ are the 2 irreducible, recurrent components. \\

$\pi_A = (2/3, 1/3)$, and $\pi_B = (8/17, 5/17, 4/17)$.
\end{minipage}
\hspace{3em}
\begin{minipage}[t]{0.5\textwidth}
$$P=\begin{array}{c|cc|cc}
& 1 & 2 & A & B \\
\hline
1 & .2 & .3 & .1 &.4 \\
2 & 0 & .5 & .2 & .3 \\
\end{array}$$

$h = R h + v$, solving gives $h = (11/40, 2/5)$.
\end{minipage} \\ \\ \\ \\

\begin{minipage}[t]{0.5\textwidth}
\setlength{\arraycolsep}{10pt}
\renewcommand{\arraystretch}{1.4}
$\lim_{n \to \infty} P^n = \begin{array}{c|c|c|c}
& T & A & B \\
\hline
T & 0 & h \pi_A & (1-h) \pi_B\\
\hline
A & 0 & \vec{1} \pi_A & 0 \\
\hline
B & 0 & 0 & \vec{1} \pi_B \\
\end{array}$
\end{minipage}
\begin{minipage}[t]{0.5\textwidth}
The limit probability of going from a transient state to a recurrent state is the probability of exiting into that recurrent state times the stationary probability of that recurrent state inside the component.
\end{minipage} \\ \\


\subsection{Exit times}
$\mathbb{E}_x[V_A]$: expected number of steps we take before hitting the exit set $A$ if we start from state $x$. \\

Let $\mathbb{E}_x[V_A] = g(x)$. $g(a) = 0$ for all $a \in A$. \\

$g(x) = \sum_y p(x, y) (1 + g(y)) = \sum_y p(x, y) + \sum_y p(x, y) g(y)$, so \boxed{g(x) = 1 + \sum_y p(x, y) g(y)}.


\subsubsection{Example}
We continue with the gambler who has a 50/50 chance of winning/losing (\ref{subsec:fair_gambler}). $S = \{0, 1, 2, 3, 4\}$, and $A = \{0, 1\}$.

\noindent
\begin{minipage}[t]{0.2\textwidth}
\begin{align*}
g(0) &= 0 \\
g(1) &= 1 + \tfrac{1}{2} g(2) \\
g(2) &= 1 + \tfrac{1}{2} g(1) + \tfrac{1}{2} g(3) \\
g(3) &= 1 + \tfrac{1}{2} g(2) \\
g(4) &= 0
\end{align*}
\end{minipage}
\hspace{1.5em}
\begin{minipage}[t]{0.40\textwidth}
$$\left(
\begin{array}{c|ccc|c}
1   & 0   & 0   & 0   & 0 \\
\hline
0  & 1  & -.5  & 0   & 0 \\
0   & -.5  & 1  & -.5  & 0 \\
0   & 0   & -.5  & 1  & 0 \\
\hline
0   & 0   & 0   & 0   & 1
\end{array}
\right)
\begin{pmatrix}
g(0) \\
g(1) \\
g(2) \\
g(3) \\
g(4)
\end{pmatrix}
=
\begin{pmatrix}
0 \\ 1 \\ 1 \\ 1 \\ 0
\end{pmatrix}$$
\end{minipage}
\hspace{2em}
\begin{minipage}[t]{0.24\textwidth}
\vspace*{1\baselineskip}
\boxed{g = \vec{1} + R g} \\

$(I-R) g = \vec{1}$ \\

$g = (I-R)^{-1} \ \vec{1}$
\end{minipage} \\ \\


\end{document}
